% AUTO-GENERATED by scripts/generate_tables.py from catalog/*.yaml -- do not edit by hand.

\begin{table*}[t]
\centering
\scriptsize
\caption{Era summary. Dates are provisional (see Sec.~\ref{sec:methodology}); representative systems are the catalog entries with the highest reading priority in each era.}
\label{tab:eras}
\begin{tabular}{p{0.11\textwidth}p{0.12\textwidth}p{0.14\textwidth}p{0.14\textwidth}p{0.14\textwidth}p{0.10\textwidth}p{0.15\textwidth}}
\toprule
\textbf{Era} & \textbf{Machine driver} & \textbf{Programming abstraction} & \textbf{Decomposition representation} & \textbf{Generated communication} & \textbf{Representative systems} & \textbf{Limitation} \\
\midrule
0: Foundations (1950--1979) & Stored-program machines, FORTRAN, pipelined vector processors, lock-step array processors, dataflow and systolic proposals & Sequential FORTRAN; vector statements; early parallel-language proposals & None for distribution; dependence graphs and loop-level analysis for vectorization & None (shared or centralised memory) & -- & No notion of data placement as a compiler problem; machines had a single address space or SIMD broadcast control \\
1: Emergence of distributed-memory compilation (1980--1987) & Multicomputers and hypercubes (Cosmic Cube, iPSC, nCUBE), systolic arrays, Connection Machine & Explicit message passing; early global-array proposals; SIMD data-parallel languages & Programmer-managed; first compiler proposals for placement & Hand-written sends/receives & -- & Programming by hand was error-prone; portability across machines was poor \\
2: Global arrays and data decomposition (1988--1992) & Commercial hypercubes and MPPs (iPSC/2, iPSC/860, nCUBE-2, CM-2) & Sequential or global-array programs with distribution annotations (SUPERB, Kali, Vienna Fortran, Fortran D) & BLOCK/CYCLIC distributions, alignment to templates/decompositions, overlap regions & Compiler-generated point-to-point messages; message vectorization; first inspector-executor runtimes & \cite{callahan1988compiling,koelbel1991compiling,hiranandani1992compiling,li1991alignment} & Regular, affine programs only; automatic distribution selection intractable in general; communication costs poorly modelled \\
3: HPF, MPI, and compiler maturity (1993--2000) & MPPs (CM-5, Paragon, SP2, T3D/T3E) and early clusters & HPF directives; MPI library calls; irregular extensions; early PGAS proposals (Co-array Fortran, Titanium) at the end of the era & HPF ALIGN/DISTRIBUTE with BLOCK, CYCLIC(k), templates and processor arrangements; INDIRECT/mapping arrays in HPF-2 and research systems & Communication sets, local address generation, redistribution, collectives; PARTI/CHAOS schedules for irregular code & \cite{zima1993compiling,benkner1999compiling,banerjee1995paradigm,loveman1993hpf} & HPF compilers were slow to mature and performance was unpredictable; explicit MPI became the dominant practice \\
4: MPI dominance and alternative programming models (2000--2010) & Commodity clusters, early multicore nodes, Cell/BE and explicitly managed memories & MPI(+OpenMP); PGAS (UPC, Co-array Fortran, Titanium); HPCS languages (Chapel, X10); hierarchical memory models (Sequoia) & Language-level distributions and locales/places; hierarchical task trees & One-sided communication, compiler-scheduled bulk transfers for memory hierarchies & \cite{tang2000tiled,chen2003upc,charles2005x10,diaconescu2005reusable} & Adoption of new languages remained limited; automatic distribution rarely used in production \\
5: Compiler/runtime co-design and accelerators (2010--2019) & GPU clusters, heterogeneous nodes, early TPU pods & Affine loop nests; task/region programs (Legion, Regent); tensor DSLs; ML graph frameworks & Polyhedral partitions; logical regions and partitions; named-dimension layouts (Mesh-TensorFlow) & Polyhedral communication-set generation with MPI; dynamic region coherence; collectives in ML frameworks & \cite{bondhugula2013compiling,kruse2014molly,bondhugula2011automatic,bauer2012legion} & Each system covered one domain; no shared ownership IR; ML partitioning still largely manual \\
6: Mesh sharding, ML systems, and spatial accelerators (2020--2026) & TPU/GPU pods, AI accelerators, wafer-scale engines, spatial dataflow chips & Sharding annotations on device meshes; tensor algebra with distribution formats; MLIR dialects; stencil DSLs & Device meshes with per-dimension sharding; distributed tensor formats; MLIR shard/mesh attributes & Collective insertion by propagation; resharding; automated search over parallelism plans & \cite{xu2021gspmd,zheng2022alpa,alabed2025partir,yadav2022distal} & Dynamic/irregular workloads, heterogeneous hierarchies, and memory-capacity constraints remain weakly handled; fragmentation across frameworks \\
\bottomrule
\end{tabular}
\end{table*}
